\section{Results}
\label{sec:results}

Our pipeline successfully identified several highly active, low-frequency features in the \swissprot dataset. \Secref{sec:results:features} details the quantitative properties of these features, and \secref{sec:results:hypotheses} presents the biological hypotheses generated by the LLM.

\subsection{Identification of Orphan Features}
\label{sec:results:features}

We identified the top 5 features with the highest maximum activation that activated in fewer than 90\% of the sequences. These features represent localized structural or functional motifs rather than global sequence properties. \Tabref{tab:feature_stats} summarizes their activation statistics.

\begin{table}[h]
\centering
\caption{Activation statistics for the top 5 identified orphan features in \esm layer 4. Frequency is defined as the fraction of positions where activation $> 0.1$.}
\begin{tabular}{@{}lccc@{}}
\toprule
\textbf{Feature Index} & \textbf{Max Activation} & \textbf{Frequency ($>0.1$)} \\
\midrule
4383 & {\bf 2.32} & 0.01 \\
9934 & 1.33 & 0.01 \\
608  & 1.30 & 0.01 \\
716  & 1.25 & 0.01 \\
989  & 1.09 & 0.02 \\
\bottomrule
\end{tabular}
\label{tab:feature_stats}
\end{table}

The most prominent feature, Index 4383, exhibited a peak activation of 2.32, which is significantly higher than the average activation across the \sae. The low frequencies (1--2\%) confirm that these features are highly specific to particular sequence contexts.

\subsection{LLM-Generated Biological Hypotheses}
\label{sec:results:hypotheses}

We utilized \gpt to interpret the sequence windows associated with the peak activations of these features. \Tabref{tab:hypotheses} presents the hypothesized motifs and the corresponding biophysical rationale.

\begin{table}[h]
\centering
\caption{LLM-generated hypotheses for selected \sae features. Sequence windows show 10 residues flanking the peak activation site (indicated in brackets).}
\resizebox{\textwidth}{!}{%
\begin{tabular}{@{}lll p{6cm}@{}}
\toprule
\textbf{Feature} & \textbf{Peak Sequence Window} & \textbf{Hypothesized Motif} & \textbf{Biophysical Rationale} \\
\midrule
4383 & \texttt{EKGKQSGIIG[A]SLDGTNPAL} & Charged Turn Motif & Presence of flanking charged residues (E, K, Q) surrounding a central alanine suggests a propensity for a flexible turn stabilized by electrostatic interactions. \\
\midrule
9934 & \texttt{DANAEIEILR[S]RMVVGKAVD} & Positively-Charged Helix Cap & The central arginine (R) and surrounding polar residues often act as helix stabilizers, specifically at the capping positions where charge neutralization occurs. \\
\midrule
716 & \texttt{EKGKQSGIIG[A]SLDGTNPAL} & Charged Beta-Turn & Centered glycine (G) is a common beta-turn initiator; surrounding charged residues provide local stabilization in non-helical regions. \\
\bottomrule
\end{tabular}
}
\label{tab:hypotheses}
\end{table}

Interestingly, Features 4383 and 716 both peaked on the same sequence segment from the dataset. While Feature 4383 was interpreted as a general "Charged Turn Motif," Feature 716 was characterized more specifically as a "Charged Beta-Turn" due to the emphasis on the glycine residue. This phenomenon, known as "feature splitting" \cite{simon2024interplm}, suggests that the \sae has learned multiple, slightly different facets of the same underlying structural concept, providing a more granular view of the motif than a single feature would allow.

\subsection{Random Baseline Comparison}
To verify the validity of the LLM's interpretations, we performed a baseline test using a random sequence fragment: \texttt{NPTGLKASLS[P]QRELRAQG}. When prompted with this sequence, the LLM hypothesized a "Proline-Induced Structural Disruption." While biophysically plausible (as prolines often break helices), the hypothesis was more generic and lacked the specific alignment with local charge distributions seen in the \sae-selected features. This indicates that while the LLM is capable of rationalizing any sequence, the \sae performs the critical role of selecting sequences that actually represent conserved, repeating motifs across the model's internal representation.
